El problema estudiado consiste en encontrar, dentro de una grilla con celdas normales, murallas y minas, 
un camino desde el origen hasta la meta que maximice el puntaje sin agotar las vidas disponibles. 
Se comparan los algoritmos de fuerza bruta, programación dinámica y 
una estrategia greedy mediante sus tiempos, memoria y calidad de solución. 
El objetivo es contrastar el comportamiento medido con las complejidades teóricas
 y caracterizar cuándo la decisión local de greedy pierde calidad o no alcanza la meta.